% Pre-experiments draft. Sections written: Abstract, Intro, Related Work, Data,
% Experimental Design, Planned Analysis. Results/Discussion/Conclusion are stubs
% to be filled in once the runs complete. Uses the official ACL submission
% template (review mode = with line ruler).

\documentclass[11pt]{article}
\usepackage[preprint]{acl}

\usepackage{times}
\usepackage{latexsym}
\usepackage[T1]{fontenc}
\usepackage[utf8]{inputenc}
\usepackage{microtype}
\usepackage{inconsolata}
\usepackage{graphicx}
\usepackage{booktabs}
\usepackage{amsmath}

\title{Does Translation Ambiguity Modulate Sycophancy in
       LLM-Based Machine Translation Evaluation?}

\author{Author One \\ McGill University \\ \texttt{author1@mail.mcgill.ca}
        \And
        Author Two \\ McGill University \\ \texttt{author2@mail.mcgill.ca}
        \And
        Author Three \\ McGill University \\ \texttt{author3@mail.mcgill.ca}}

\begin{document}
\maketitle

\begin{abstract}
LLMs are now widely used to judge machine translation (MT) quality, but
recent work outside MT shows these judges can be talked into revising their
scores when a user pushes back. We ask whether this effect depends on how
ambiguous the translation is in the first place: do LLM judges cave more on
segments where human raters themselves disagreed? We measure ambiguity per
segment as the standard deviation of standardised human direct-assessment
(DA) scores from WMT Quality Estimation, and draw a 594-segment sample
balanced over a $3 \times 3$ gold-quality $\times$ ambiguity grid so the two
axes are decorrelated. The challenge protocol uses two controls. The first
isolates random reconsideration noise; the second isolates whether the
critique itself was informative. What is left is attributable to social
pressure. Results to follow.
\end{abstract}

\section{Introduction}

LLMs are increasingly used to judge MT
quality~\cite{kocmi-federmann-2023-gemba,wmt2025metrics}. A separate line
of work has shown that LLM evaluators can be \emph{sycophantic} under user
rebuttal: they revise their judgements toward whatever the user suggests,
even when the user is
wrong~\cite{sharma-sycophancy-2023,kim-khashabi-2025,hwang-grader-2025}.
Whether this matters for MT, and when, is open.

We hypothesise that sycophancy is stronger on translations where human
raters themselves disagreed about quality. When a translation is
borderline, the judge has less internal conviction to defend, and is
easier to push around. Concretely we test three claims:
\textbf{H1} LLM MT judges yield to user pressure: when challenged, they
revise their scores in the direction the user suggests, even when the
user is wrong;
\textbf{H2} (central) the size of that revision grows with per-segment
human disagreement;
\textbf{H3} rubric prompts that require explicit reasoning resist
sycophancy better than scalar score-only prompts.

We test this jointly: LLM-judge sycophancy under user challenge against
per-segment human-rater disagreement. We use the English\,$\to$\,Hindi
split of WMT QE~2023, because the test requires per-rater scalar DA
scores rather than aggregated means or categorical error spans
(Section~\ref{sec:dataset}).

\section{Related Work}

\paragraph{LLM judges for MT.}
GEMBA~\cite{kocmi-federmann-2023-gemba} showed that zero-shot prompted LLM
judges are competitive with learned MT metrics;
GEMBA-MQM~\cite{kocmi-federmann-2023-gemba-mqm} extends this to MQM-style
error spans, and RUBRIC-MQM~\cite{kim-rubric-mqm-2025} meta-evaluates MT
LLM judges. The WMT25 metrics shared task~\cite{wmt2025metrics} reports
that strong LLM judges top system-level rankings, while segment-level
prediction is harder. None of this work probes what happens when a user
pushes back.

\paragraph{Sycophancy and persuasion.}
Sharma et al.~\cite{sharma-sycophancy-2023} framed the basic phenomenon:
LLMs systematically agree with stated user beliefs.
SycEval~\cite{syceval-2025} introduced the \emph{progressive} vs.\
\emph{regressive} distinction we adopt in Section~\ref{sec:outcomes}
(regressive: the revision moves further from gold; progressive: closer).
Kim and Khashabi~\cite{kim-khashabi-2025} are the closest conceptual
neighbour: they show LLM evaluators flip more readily under follow-up
rebuttal than under neutral evaluation, especially for casually-worded
challenges. Neither study is MT-specific, and neither conditions on
human-rater disagreement.

\paragraph{Gap.}
The streams above have been studied separately. The joint test ---
\emph{does MT judge sycophancy concentrate on the segments humans found
hardest to score?} --- is, to our knowledge, untested.

\section{Data}

\subsection{Dataset requirement and choice}
\label{sec:dataset}

We need \emph{per-segment, per-rater scalar quality scores}, since
ambiguity is the SD across raters for a single segment. This rules out
two common MT annotation regimes. MQM
annotations~\cite{kocmi-federmann-2023-gemba-mqm} record categorical
error spans, so inter-rater dispersion has no natural continuous
definition. And several recent DA releases ship only aggregated
$z$-means after pooling raters, which discards the dispersion we need.

WMT QE~2023~\cite{wmt-qe-2023} fits cleanly: it releases sentence-level
DA from four expert raters per segment, with per-rater $z$-standardisation
applied by WMT before release. We use its English\,$\to$\,Hindi train and
development splits (8{,}000 segments combined).

\subsection{Gold and ambiguity}

For each segment $i$ with raters $r = 1, \ldots, k$ ($k=4$), let $z_{i,r}$
be the released $z$-standardised DA score. We define
\begin{align*}
g_i &= \frac{1}{k}\sum_{r=1}^{k} z_{i,r}                & \text{(gold)} \\
A_i &= \mathrm{SD}\bigl(z_{i,1}, \ldots, z_{i,k}\bigr)  & \text{(ambiguity)}.
\end{align*}
$g_i$ is our best estimate of segment quality and $A_i$ measures how much
the four raters disagreed.

\subsection{Cross-stratified sample}
\label{sec:sample}

The full en$\to$hi data has $\mathrm{Spearman}(g, A) = -0.31$:
lower-quality translations are systematically more disagreed upon.
Stratifying only on $A$ would bias the high-ambiguity tertile toward bad
translations and let a reviewer attribute any H2 effect to translation
badness rather than ambiguity. We therefore cross-stratify on a
$3 \times 3$ grid of (gold-tertile $\times$ ambiguity-tertile), drawing
66 segments from each of the 9 cells ($n=594$, seed 42), with cell
boundaries set by terciles of $g$ and $A$ in the full table. Residual
$\mathrm{Spearman}(g, A)$ in the sample is $-0.13$.

\section{Experimental Design}

\subsection{Judges and prompts}

We use two frontier LLM judges from independent providers, with temperature
fixed at 0 and exact snapshot IDs pinned in the released code. Each judge
is queried under two prompt families. The \textbf{scalar} prompt asks only
for a 0--100 score (WMT25 template). The \textbf{rubric} prompt asks the
model to first identify any major errors with a brief reason, then emit a
JSON object \texttt{\{score\_0\_100, major\_issue, brief\_reason\}}. The
contrast tests H3.

\subsection{Conditions: one challenge and two controls}

For every (segment, judge, prompt-family) trio we run four queries, using
a single \textbf{casual-doubt} challenge style.

\begin{enumerate}
\item \textbf{Baseline.} The judge sees source and translation, emits $s_0$.
\item \textbf{Challenge (same chat).} The user follows up with a casual
objection in the \emph{gold-inconsistent} direction: if humans agreed the
translation is good, the user pushes for a lower score, and vice versa.
The judge emits $s_c$. Because we know the gold direction, any compliance
is sycophantic rather than corrective.
\item \textbf{Neutral control (same chat).} Same context, but the follow-up
is a neutral ``Could you take another look at your score?'' with no opinion.
Emits $s_{\mathrm{neut}}$, controlling for second-look drift, in the spirit
of the neutral-reprompt baselines used in prior sycophancy
work~\cite{sharma-sycophancy-2023}.
\item \textbf{Simultaneous control (fresh chat).} A new conversation in
which the objection appears alongside source and translation from the start,
with no prior baseline to defend. Emits $s_{\mathrm{sim}}$, controlling for
the critique's informational content independent of social pressure. This
is the same mechanism as the LLM-as-Judge contrast
in~\citet{kim-khashabi-2025}, ported here from MCQ answers to MT scores.
\end{enumerate}

\subsection{Outcome measures}
\label{sec:outcomes}

For challenged score $s_c$, baseline $s_0$, and gold $g$, we record:
\begin{itemize}
\item \emph{Signed shift:} $s_c - s_0$.
\item \emph{Directional compliance:} $|s_c - s_0| \geq 5$ in the
user-requested direction.
\item \emph{Regressive sycophancy:} compliance such that
$|s_c - g| > |s_0 - g|$.
\item \emph{Progressive correction:} compliance such that
$|s_c - g| < |s_0 - g|$.
\end{itemize}
The progressive/regressive split is adapted from SycEval~\cite{syceval-2025},
which introduced it for direct answers.

\section{Planned Analysis}

\begin{description}
\item[H1 (sycophancy exists).] Paired Wilcoxon on
$|s_c - g| - |s_0 - g|$ per (judge $\times$ prompt), with the neutral
control as the within-segment null.
\item[H2 (ambiguity moderates).] Spearman $\rho$ between $A$ and
$|s_c - s_0|$, with bootstrap 95\% CI ($n=10{,}000$).
\item[H3 (rubric resists).] Paired Wilcoxon on directional-compliance
rate, scalar vs.\ rubric.
\end{description}

As a robustness check we fit a mixed-effects regression
\[
\mathrm{shift} \sim A_z + g_z + \mathrm{prompt} + \mathrm{judge}
                + (1 \mid \mathrm{segment}),
\]
with $g_z$ as a covariate to control for the residual quality--ambiguity
correlation (Section~\ref{sec:sample}).

\section{Results}

\section{Discussion}

\section{Conclusion}

\section*{Statement of Contributions}

\bibliographystyle{acl_natbib}
\bibliography{references}

\end{document}
